\documentclass[UTF8,zihao=-4]{ctexart}
\usepackage[a4paper,margin=2.1cm]{geometry}
\usepackage{hyperref}
\usepackage{array}
\usepackage{longtable}
\usepackage{verbatim}
\hypersetup{hidelinks}
\setlength{\parindent}{0pt}
\setlength{\parskip}{0.55em}
\emergencystretch=4em
\sloppy
\title{03-语义分析与中间代码复习讲义}
\author{}
\date{}
\begin{document}
\maketitle
\setcounter{tocdepth}{1}
\tableofcontents
\newpage
\section{03 语义分析与中间代码复习讲义}


对应资料：


\begin{itemize}
\item \texttt{Slides/Slides/Ch6-1-Intermediate Code Generation.pdf}
\item \texttt{Slides/Slides/Ch6-2-Translation into SSA.pdf}
\item \texttt{Slides/Slides/Slides-Ant/7-symtable-handout.pdf}
\item \texttt{Slides/Slides/Slides-Ant/8-type-systems-out.pdf}
\item \texttt{Slides/Slides/Slides-Ant/9-semantics-ag-handout.pdf}
\item \texttt{Slides/Slides/Slides-Ant/10-llvm-handout.pdf}
\item \texttt{Slides/Slides/Slides-Ant/11-ir-expr-control-handout.pdf}
\item \texttt{Slides/Slides/handout/3 IR SSA.pdf}
\item \texttt{Review/compile.pdf}
\item \texttt{Exam/2025.pdf} 支配与 SSA 相关题
\item \texttt{Exam/过时/2022-exam-A.pdf} 题目 5、6
\item \texttt{Exam/过时/2023-exam-A.pdf} 题目 5
\item \texttt{Exam/过时/2021-exam-A.pdf} 题目 4、5
\end{itemize}


\subsection{一、考试目标}


本章高频题：


\begin{quote}
\footnotesize
\begin{verbatim}
语法制导定义：设计属性和语义规则
中间代码：写三地址码
回填：truelist、falselist、nextlist
SSA：支配、支配边界、phi
\end{verbatim}
\end{quote}


零基础顺序：


\begin{quote}
\footnotesize
\begin{verbatim}
先会写属性表，再会写回填模板，最后背 SSA 定义。
\end{verbatim}
\end{quote}


\subsection{二、语义分析}


语义分析检查语法树中的“含义”：


\begin{itemize}
\item 类型检查。
\item 作用域检查。
\item 声明与使用匹配。
\item 函数调用参数个数和类型检查。
\end{itemize}


符号表常用操作：


\begin{quote}
\footnotesize
\begin{verbatim}
enterScope()
exitScope()
define(name, type, attributes)
lookup(name)
\end{verbatim}
\end{quote}


\subsection{三、语法制导定义 SDD}


SDD 把属性和语义规则挂在产生式上。


两类属性：


\begin{quote}
\footnotesize
\begin{verbatim}
综合属性：由子节点算父节点。
继承属性：由父节点或左侧兄弟节点传给当前节点。
\end{verbatim}
\end{quote}


答题格式：


\begin{quote}
\footnotesize
\begin{verbatim}
属性定义：
  X.val : 综合属性，表示数值
  X.len : 综合属性，表示长度

产生式              语义规则
S -> L . L          S.val = L1.val + L2.val / 2^(L2.len)
L -> L B            L.val = L1.val * 2 + B.val; L.len = L1.len + 1
B -> 0              B.val = 0
B -> 1              B.val = 1
\end{verbatim}
\end{quote}


\subsection{四、SDD 真题模板}


二进制小数转十进制：


\begin{quote}
\footnotesize
\begin{verbatim}
S -> L1 . L2   S.val = L1.val + L2.val / 2^(L2.len)
S -> L         S.val = L.val
L -> L1 B      L.val = L1.val * 2 + B.val; L.len = L1.len + 1
L -> B         L.val = B.val; L.len = 1
B -> 0         B.val = 0
B -> 1         B.val = 1
\end{verbatim}
\end{quote}


统计声明个数：


\begin{quote}
\footnotesize
\begin{verbatim}
P -> D              P.count = D.count
D -> D1 ; D2        D.count = D1.count + D2.count
D -> T id           D.count = 1
D -> func id {D1}   D.count = 1 + D1.count
T -> int            不改变 count
\end{verbatim}
\end{quote}


打印变量嵌套深度：


\begin{quote}
\footnotesize
\begin{verbatim}
P -> D              D.depth = 0
D -> D1 ; D2        D1.depth = D.depth; D2.depth = D.depth
D -> T id           print(id, D.depth)
D -> func id {D1}   D1.depth = D.depth + 1
T -> int            不改变 depth
\end{verbatim}
\end{quote}


C 声明英文说明模板：


\begin{quote}
\footnotesize
\begin{verbatim}
D -> id       D.name = id.lexeme; D.desc = ""
D -> D1 ()    D.name = D1.name; D.desc = D1.desc + " function returning"
D -> D1 []    D.name = D1.name; D.desc = D1.desc + " array"
D -> * D1     D.name = D1.name; D.desc = D1.desc + " pointer to"
D -> ( D1 )   D.name = D1.name; D.desc = D1.desc
输出：declare D.name as D.desc
\end{verbatim}
\end{quote}


\subsection{五、三地址码 TAC}


三地址码每条指令最多包含三个地址。


常见形式：


\begin{quote}
\footnotesize
\begin{verbatim}
x = y op z
x = op y
x = y
goto L
if x relop y goto L1
if x relop y goto L1 else goto L2
param x
call f, n
return x
\end{verbatim}
\end{quote}


数组访问在考试中可简化为表达式：


\begin{quote}
\footnotesize
\begin{verbatim}
if dict[mid] < key goto L
\end{verbatim}
\end{quote}


\subsection{六、回填}


回填用于控制流翻译。先把跳转目标留空，之后统一填标签。


核心属性：


\begin{quote}
\footnotesize
\begin{verbatim}
B.truelist  : B 为真时要跳转的指令编号集合
B.falselist : B 为假时要跳转的指令编号集合
S.nextlist  : S 执行完后要跳转的指令编号集合
\end{verbatim}
\end{quote}


核心函数：


\begin{quote}
\footnotesize
\begin{verbatim}
makelist(i)    创建只含 i 的链表
merge(p, q)    合并链表
backpatch(p,L) 把 p 中所有空目标填成 L
nextquad       下一条指令编号
\end{verbatim}
\end{quote}


布尔表达式：


\begin{quote}
\footnotesize
\begin{verbatim}
B -> E1 relop E2
  B.truelist = makelist(nextquad)
  emit("if E1 relop E2 goto _")
  B.falselist = makelist(nextquad)
  emit("goto _")
\end{verbatim}
\end{quote}


\subsection{七、while 回填模板}


语句：


\begin{quote}
\footnotesize
\begin{verbatim}
while B do S1
\end{verbatim}
\end{quote}


模板：


\begin{quote}
\footnotesize
\begin{verbatim}
M1.quad = nextquad          循环条件入口
翻译 B
M2.quad = nextquad          循环体入口
backpatch(B.truelist, M2.quad)
翻译 S1
emit("goto M1.quad")
backpatch(S1.nextlist, M1.quad)
S.nextlist = B.falselist
\end{verbatim}
\end{quote}


二分查找与 GCD 都按这个模板：


\begin{quote}
\footnotesize
\begin{verbatim}
循环条件为真 -> 进入循环体
循环体内部 if/else 继续回填
循环体结束 -> goto 条件入口
循环条件为假 -> 跳出 while
\end{verbatim}
\end{quote}


\subsection{八、回填最小例子：GCD}


源程序：


\begin{quote}
\footnotesize
\begin{verbatim}
while a != b do
  if a > b then
    a = a - b
  else
    b = b - a
\end{verbatim}
\end{quote}


一种编号写法：


\begin{quote}
\footnotesize
\begin{verbatim}
100: if a != b goto 102
101: goto 108
102: if a > b goto 104
103: goto 106
104: a = a - b
105: goto 100
106: b = b - a
107: goto 100
108: ...
\end{verbatim}
\end{quote}


说明：


\begin{quote}
\footnotesize
\begin{verbatim}
B.truelist 回填到 102。
B.falselist 成为 while 的 nextlist，回填到 108。
if 的真分支到 104，假分支到 106。
两个分支结束都回到 100。
\end{verbatim}
\end{quote}


\subsection{九、SSA}


SSA：Static Single Assignment，每个变量只被定义一次。


普通代码：


\begin{quote}
\footnotesize
\begin{verbatim}
x = 1
x = x + 1
\end{verbatim}
\end{quote}


SSA：


\begin{quote}
\footnotesize
\begin{verbatim}
x1 = 1
x2 = x1 + 1
\end{verbatim}
\end{quote}


控制流汇合处使用 phi：


\begin{quote}
\footnotesize
\begin{verbatim}
if ...
  x1 = 1
else
  x2 = 2
x3 = phi(x1, x2)
\end{verbatim}
\end{quote}


phi 含义：


\begin{quote}
\footnotesize
\begin{verbatim}
根据控制流来自哪个前驱块，选择对应版本。
\end{verbatim}
\end{quote}


\subsection{十、支配关系}


定义：


\begin{quote}
\footnotesize
\begin{verbatim}
Dom(n)：从入口到 n 的每条路径都经过的节点集合。
idom(n)：n 的最近严格支配者。
Dom Frontier(n)：n 支配某个前驱，但不严格支配该节点的位置。
PostDom(n)：从 n 到出口的每条路径都经过的节点集合。
\end{verbatim}
\end{quote}


SSA 中 phi 放置：


\begin{quote}
\footnotesize
\begin{verbatim}
变量 x 在多个基本块中被定义。
在这些定义块的支配边界处放置 x 的 phi。
\end{verbatim}
\end{quote}


支配树：


\begin{quote}
\footnotesize
\begin{verbatim}
除入口外，每个节点有唯一 idom。
idom 关系无环。
因此支配关系图是以入口为根的树；多入口时是森林。
\end{verbatim}
\end{quote}


\subsection{十一、自测题}


题 1：综合属性和继承属性区别？


答案：


\begin{quote}
\footnotesize
\begin{verbatim}
综合属性由子节点向父节点传递。
继承属性由父节点或左侧兄弟节点向当前节点传递。
\end{verbatim}
\end{quote}


题 2：回填三件套是什么？


答案：


\begin{quote}
\footnotesize
\begin{verbatim}
makelist、merge、backpatch。
\end{verbatim}
\end{quote}


题 3：while 的假出口是什么？


答案：


\begin{quote}
\footnotesize
\begin{verbatim}
条件 B 的 falselist。
\end{verbatim}
\end{quote}


题 4：SSA 中 phi 放在哪里？


答案：


\begin{quote}
\footnotesize
\begin{verbatim}
放在变量定义块的支配边界处。
\end{verbatim}
\end{quote}

\end{document}
